\documentclass[10pt,a4paper]{amsart}
\usepackage[margin=19mm]{geometry}
\usepackage{amsmath,amssymb,mathtools}
\usepackage[scr=rsfs,cal=euler]{mathalfa}
\usepackage{xfrac}
\usepackage[unicode,hidelinks,bookmarks=false]{hyperref}
\newcommand{\ie}{i.e.\ }
\newcommand{\cf}{cf.\ }
\newcommand{\resp}{respectively\ }
\newcommand{\Subsection}{\subsection}
\providecommand{\nobreakdash}{\nobreak\hbox{-}}
\setlength{\emergencystretch}{2em}
\multlinegap0pt
\usepackage{amssymb}
\newtheorem{lemma}{Lemma}
\newtheorem{theorem}{Theorem}
\newcommand{\Z}{\mathbb{Z}}
\newcommand{\killop}{\operatorname{kill}}
\title{Exact values and exact upper bounds for families of integers with arithmetic progression intersections\\ \large (Erd\H{o}s Problem \#272)}
\author{Zhanfu Yang}
\date{Source: arXiv:2607.23004v1 (CC BY 4.0)}
\begin{document}
\maketitle

\noindent\emph{Source and licence.} Exact values and exact upper bounds for families of integers with arithmetic progression intersections\\ \large (Erd\H{o}s Problem \#272). Version arXiv:2607.23004v1 (CC BY 4.0), distributed by arXiv under the Creative Commons Attribution 4.0 International licence, http://creativecommons.org/licenses/by/4.0/, as recorded in the arXiv licence metadata for this exact version and shown on the abstract page https://arxiv.org/abs/2607.23004v1. Full licence notice: \texttt{licenses/arxiv-2607.23004-CC-BY-4.0.txt}.\par\smallskip

\noindent\textit{Editorial scope.} Counted unit: the article's theorem thm:private (source0.tex lines 197--199). The article's complete original proof of it is delivered with it (source0.tex lines 201--213, one \texttt{proof} environment of 13 lines). No mathematics is written, repaired or re-derived: the statement and the proof are the article's own lines, in the article's own notation, and no retained line is substituted or re-flowed.  Retained uncounted as the source-local context the proof needs: lines 119--121.  Nothing was omitted from any retained span, so the package carries no omission record.  No retained line contains a citation, so the wrapper carries no bibliography and no bibliography entry.  No retained line contains a figure environment, an image-inclusion command or drawing code, so no image is delivered and the package declares no asset dependency.  Editorial formatting only, in addition: an AMS \texttt{amsart} preamble that restates the article's own declarations and macros used by the retained lines, namely the environments \texttt{lemma}, \texttt{theorem} on separate counters, together with the standard packages those lines need.  The retained environments print in this document as Lemma 1, Theorem 1 in order of appearance, whereas the article prints the same items with the numbering of its own full document.  The complete native source of the article as distributed in the arXiv e-print for this exact version is bundled unchanged as \texttt{sources/205992/source0.tex}.
\begin{lemma}[Triangle-freeness]\label{lem:tri}
$R$ contains no triangle. Consequently every $z$ with $|z|\ge 4$ contains a bad pair, so a single member has $|\{z\}|-\killop(\{z\})\le 0$.
\end{lemma}

\begin{theorem}[Private-pair theorem]\label{thm:private}
Every crooked member contains a private bad pair.
\end{theorem}

\begin{proof}
Since $|z\setminus\{0\}|\ge 3$ and the relation $R$ is triangle-free (Lemma~\ref{lem:tri}), $z$ contains a bad pair. Assume for contradiction that every bad pair of $z$ is spanned in $z$; we show $z$ is an AP.

Let $\delta_0$ be the least positive integer that spans some bad pair of $z$, say $p_0=\{\delta_0 a,\delta_0 b\}$ with $\gcd$-condition $\delta_0\mid\gcd$, and let $P_0\subseteq z$ be the corresponding progression: all multiples of $\delta_0$ in $I_0=[\min(0,\delta_0a,\delta_0b),\max(0,\delta_0a,\delta_0b)]$. Since $\{a,b\}$ is not an $R$-pair we cannot have $|a|=|b|=1$, so $\max(|a|,|b|)\ge2$; hence $P_0$ contains $2t\delta_0$ for some sign $t\in\{+,-\}$, and $P_0$ contains $s\delta_0$ for some sign $s$.

\emph{Step 1: $z\subseteq\delta_0\mathbb{Z}$.} Let $w\in z$ with $\delta_0\nmid w$. The $R$-partners of $s\delta_0$ are $-s\delta_0$, $2s\delta_0$ and $s\delta_0/2$; the first two are multiples of $\delta_0$. If $w\neq s\delta_0/2$, the pair $\{s\delta_0,w\}$ is therefore bad, and $g:=\gcd(\delta_0,|w|)$ is a proper divisor of $\delta_0$; by assumption this pair is spanned at some $\delta\mid g<\delta_0$, contradicting the minimality of $\delta_0$. If $w=s\delta_0/2$, consider instead the pair $\{s\delta_0/2,\,2t\delta_0\}$: the $R$-partners of $s\delta_0/2$ are $-s\delta_0/2$, $s\delta_0$ and $s\delta_0/4$, none of which equals $\pm2\delta_0$, so the pair is bad, with $\gcd$ equal to $\delta_0/2<\delta_0$ --- the same contradiction. Hence $z\subseteq\delta_0\mathbb{Z}$.

\emph{Step 2: $z$ is an AP.} Badness, $R$-edges and spanning are invariant under $x\mapsto x/\delta_0$ on $\delta_0\mathbb{Z}$, so we may assume $\delta_0=1$; then $P_0$ is an integer interval around $0$ of length at least $3$, so $1\in z$ or $-1\in z$.

Suppose first $1\in z$. For $w\in z$ with $w\ge3$: the pair $\{1,w\}$ is bad ($w\notin\{-1,2,\tfrac12\}$) with $\gcd$ equal to $1$, so its only possible spanning step is $\delta=1$, forcing $[0,w]\cap\Z\subseteq z$. For $w\in z$ with $w\le-2$: likewise $\{1,w\}$ is bad and $[w,1]\cap\Z\subseteq z$. The remaining elements $w\in\{-1,2\}$ impose nothing, but $[-1,0]$ and $[0,2]$ lie in $z$ automatically whenever those elements are present. Hence $z=[\min z,\max z]\cap\Z$, an AP.

If instead only $-1\in z$: for any $w\in z$ with $w\ge 2$ the pair $\{-1,w\}$ is bad ($w\notin\{1,-2,-\tfrac12\}$) with $\gcd$ equal to $1$, forcing $[-1,w]\cap\Z\subseteq z$ and in particular $1\in z$, returning us to the previous case; if $\max z\le 1$ then either $1\in z$ (previous case) or $\max z=0$, and the mirror argument with pairs $\{-1,w\}$, $w\le-3$, gives $z=[\min z,0]\cap\Z$. In every case $z$ is an AP, contradicting crookedness.
\end{proof}

\end{document}
